Tabellen mit Zahlenmaterial sollen so gestaltet werden, dass Zahlenangaben
einer Spalte am Dezimalkomma ausgerichtet sind. Das Paket~\ctanref{siunitx}
stellt hierfür den Spaltentyp~\texttt{S} bereit. Zahlen können einfach als
Dezimalzahlen eingeben werden. Das Paket analyisiert die Eingaben selbständig
und nimmt die korrekte Ausrichtung vor. Tabelle~\ref{tab:latex:tab:numcol}
zeigt ein entsprechendes Anwendungsbeispiel und
Listing~\ref{lst:latex:tab:numcol} die zugehörigen \LaTeX-Code.

\begin{table}
\centering
\langinput{code/latex_tab_numcol}
\caption{Bahnelemente der Planeten zur Epoche J2000.0}
\label{tab:latex:tab:numcol}
\end{table}

\lstinputlisting[language={[LaTeX]{TeX}},style=block,float,caption={Tabelle mit ausgerichteten Zahlenspalten},label={lst:latex:tab:numcol}]{code/latex_tab_numcol.de.tex}

Das Paket~\ctanpkg{siunitx} ist sehr flexibel konfigurierbar, was die
Ausrichtung der Zelleninhalte und die Darstellung von Zahlenwerten angeht.
Sollten Sie innerhalb Ihrer Arbeit mit Zahlenkolonnen hantieren, lohnt sich ein
Blick in das Pakethandbuch auf jeden Fall. Mit \ctanref{dcolumn} steht ein
älteres Paket zur Ausrichtung von Zahlen in Tabellen zur Verfügung.
